% ===========================================================================
% Commandes, styles et environnements réutilisables.
% Définis une seule fois ici, jamais redéfinis dans un chapitre.
% ===========================================================================

% ─── Bascules de partie du document ───────────────────────────────────────
% La classe `report` ne fournit pas \frontmatter / \mainmatter : on encapsule
% le comportement pour qu'un gabarit officiel qui les fournit puisse être
% substitué sans toucher à main.tex.
\newcommand{\piqbitFrontmatter}{%
  \pagenumbering{roman}%
  \pagestyle{plain}%
  \onehalfspacing%
}
\newcommand{\piqbitMainmatter}{%
  \cleardoublepage%
  \pagenumbering{arabic}%
  \pagestyle{fancy}%
  \fancyhf{}%
  \fancyhead[L]{\nouppercase{\leftmark}}%
  \fancyhead[R]{\thepage}%
  \renewcommand{\headrulewidth}{0.4pt}%
}

% ─── Ouverture et clôture de chapitre ─────────────────────────────────────
% Chaque chapitre s'ouvre sur son objectif et se ferme sur une synthèse :
% règle de `pfe_latex_conventions.md`, rendue mécanique ici.
\newcommand{\chapterintro}[1]{\noindent\textit{#1}\par\bigskip}
\newcommand{\chapteroutro}[1]{\par\bigskip\noindent\textit{#1}\par}

% ─── Marqueur de fait non résolu ──────────────────────────────────────────
% Rend les TODO VISIBLES dans le PDF de travail. À supprimer, avec tous les
% \todoitem restants, avant la remise : un TODO visible est une information
% manquante assumée, pas un oubli.
\newcommand{\todoitem}[1]{\textbf{[TODO : #1]}}

% ─── Figure de diagramme ──────────────────────────────────────────────────
% Enveloppe unique pour les neuf vues TikZ de report/diagrams/.
%   #1 nom du fichier (sans extension)  #2 légende  #3 étiquette
% `max width=\textwidth` ne réduit QUE si la figure dépasse la justification :
% une vue déjà à la bonne taille garde sa taille de police nominale. Sans ce
% filet, un ajout de nœud dans un diagramme déborde silencieusement en marge.
%
% `max totalheight` joue le même rôle dans l'autre dimension, qui manquait :
% une figure plus HAUTE que la page ne débordait pas en marge, elle devenait
% un flottant que LaTeX ne pouvait placer nulle part (« Float too large for
% page »), et il la poussait en fin de chapitre. 0.82 et non 1.0 parce que le
% flottant doit aussi loger sa légende et les blancs réglementaires au-dessus
% et au-dessous.
\newcommand{\pqfigure}[3]{%
  \begin{figure}[htbp]%
    \centering%
    \begin{adjustbox}{max width=\textwidth, max totalheight=0.82\textheight}%
      \input{diagrams/#1}%
    \end{adjustbox}%
    \caption{#2}%
    \label{#3}%
  \end{figure}%
}

% Capture d'écran. Même enveloppe que \pqfigure, mais l'actif est une image :
%   #1 nom du fichier (sans extension)  #2 légende  #3 étiquette
%
% Tant que le fichier n'existe pas, un CARTOUCHE NOMMÉ prend sa place et le
% rapport compile quand même. C'est délibéré : la structure doit être relisible
% avant que les 28 captures soient prises, et un emplacement vide doit se voir
% à la relecture plutôt que d'interrompre la compilation.
% Dès que `assets/screenshots/<nom>.png` existe, l'image le remplace sans
% aucune modification du texte.
\newcommand{\pqscreenshot}[3]{%
  \begin{figure}[htbp]%
    \centering%
    \IfFileExists{assets/screenshots/#1.png}{%
      \fboxsep=0pt\fboxrule=0.4pt%
      \fcolorbox{piqbitExt}{white}{%
        \includegraphics[width=0.94\textwidth,keepaspectratio]{#1}%
      }%
    }{%
      % \detokenize : les noms de capture contiennent des underscores, que
      % LaTeX prendrait autrement pour un passage en mode mathématique.
      \fboxsep=0pt\fboxrule=0.8pt%
      \fcolorbox{piqbitWarn}{piqbitWarnLite}{%
        \parbox[c][42mm][c]{0.94\textwidth}{\centering
          {\footnotesize\bfseries\color{piqbitWarn} CAPTURE À PRODUIRE}\par\smallskip
          {\ttfamily\small\color{piqbitWarn} \detokenize{#1}.png}\par\smallskip
          {\scriptsize\color{piqbitWarn} Mode opératoire :
            {\ttfamily \detokenize{report/06_screenshot_plan.md}}}%
        }%
      }%
    }%
    \caption{#2}%
    \label{#3}%
  \end{figure}%
}

% Variante paysage, pour les vues qui ne tiennent pas en portrait sans
% devenir illisibles (modèle de données).
\newcommand{\pqfigurewide}[3]{%
  \begin{sidewaysfigure}%
    \centering%
    \begin{adjustbox}{max width=\textheight}%
      \input{diagrams/#1}%
    \end{adjustbox}%
    \caption{#2}%
    \label{#3}%
  \end{sidewaysfigure}%
}

% ─── Étiquette d'état de preuve ───────────────────────────────────────────
% Vocabulaire imposé par validation_evidence_inventory.md. L'utiliser dans les
% tableaux de validation interdit d'écrire « supposé réussi ».
\newcommand{\evstate}[1]{\textsc{#1}}

% ─── Colonne de tableau à largeur fixe, alignée en haut ───────────────────
% \RaggedRight (ragged2e) et non \raggedright : ce dernier pose un ressort
% infini à droite, donc TeX ne césure JAMAIS et un mot trop long déborde en
% marge sans que rien ne puisse l'en empêcher. Avec un ressort fini, la
% césure redevient possible et les colonnes étroites se tiennent seules —
% sans qu'il faille régler une largeur à la main à chaque mot un peu long.
% Le `\hspace{0pt}` n'est pas décoratif : TeX ne césure QUE les mots précédés
% d'un ressort, et le premier mot d'un paragraphe suit la boîte d'indentation,
% jamais du ressort. Dans une cellule qui ne contient qu'un mot long —
% « Confidentialité », « Synchronisation » — ce mot est donc incésurable et
% déborde, même avec \RaggedRight. Le ressort nul rétablit la césure.
% `\hbadness=10000` n'est pas là pour faire taire un défaut, mais parce que la
% mesure ne s'applique pas : « underfull » signifie « ligne plus courte que la
% justification », ce qui est la DÉFINITION même d'une colonne au fer à
% gauche. TeX note la badness par rapport à un texte justifié ; dans une
% colonne ragged, l'avertissement est structurellement inévitable et ne
% signale rien. Le laisser actif noierait un vrai problème le jour où il
% surviendra. Le réglage est local à la cellule.
\newcolumntype{L}[1]{>{\RaggedRight\arraybackslash\hbadness=10000\relax\hspace{0pt}}p{#1}}

% ─── Colonne X de tabularx : même traitement ──────────────────────────────
% Par défaut `X` est JUSTIFIÉE. Dans une colonne de 4 cm, la justification
% force TeX à étirer les espaces jusqu'à créer des lignes visiblement lâches
% (c'est ce que signalaient les derniers « Underfull »), et elle jurait avec
% les colonnes L voisines, elles au fer à gauche : un même tableau avait deux
% alignements. On aligne donc X sur L.
% Conséquence visible : le bord droit des colonnes X devient irrégulier. Pour
% revenir en arrière, supprimer ces deux lignes — rien d'autre n'en dépend.
\renewcommand{\tabularxcolumn}[1]{%
  >{\RaggedRight\arraybackslash\hbadness=10000\relax\hspace{0pt}}p{#1}}

% ─── Style d'extrait de code ──────────────────────────────────────────────
% Aucun secret, aucune adresse interne, aucun chemin de poste de travail ne
% doit apparaître dans un listing.
\lstdefinestyle{piqbit}{
  basicstyle=\ttfamily\footnotesize,
  keywordstyle=\bfseries,
  commentstyle=\itshape,
  breaklines=true,
  showstringspaces=false,
  frame=single,
  numbers=left,
  numberstyle=\tiny,
  captionpos=b,
  tabsize=2
}
\lstset{style=piqbit}
